\vspace{0.8cm}

\begin{center}
{\Large \textbf{Abstract}}
\end{center}

\noindent
TunSociety is a full-stack web application designed for small Tunisian communities such as student clubs, associations, and local initiatives. The objective is not to reproduce a public social network, but to provide an internal workspace where members can identify each other, publish updates, discuss posts, create controlled connections, exchange private messages, and follow activity through notifications. The project also gives responsible actors a supervision layer: moderators can inspect risky content, while administrators can manage roles, warnings, account freezes, appeals, user risk information, and audit logs. The solution was implemented with Angular, ASP.NET Core, Entity Framework Core, MySQL, JWT authentication, BCrypt password hashing, and a local AI-assisted moderation service based on Ollama and Gemma. The main contribution of the project is the combination of social interaction and community governance in one platform adapted to the needs of organized groups.

\vspace{0.6cm}

\begin{center}
{\Large \textbf{R\'esum\'e}}
\end{center}

\noindent
TunSociety est une application web full-stack destin\'ee aux petites communaut\'es tunisiennes, notamment les clubs \'etudiants, les associations et les initiatives locales. L'objectif n'est pas de copier un r\'eseau social public, mais de proposer un espace interne o\`u les membres peuvent s'identifier, publier, commenter, r\'eagir, cr\'eer des liens contr\^ol\'es, \'echanger des messages priv\'es et suivre l'activit\'e gr\^ace aux notifications. La plateforme ajoute aussi une couche de supervision: les mod\'erateurs peuvent consulter les contenus \`a risque, tandis que les administrateurs peuvent g\'erer les r\^oles, les avertissements, les gels de comptes, les appels, les indicateurs de risque et les journaux d'audit. La solution utilise Angular, ASP.NET Core, Entity Framework Core, MySQL, JWT, BCrypt et une mod\'eration locale assist\'ee par IA avec Ollama et Gemma. L'apport principal du projet est l'association entre interaction sociale et gouvernance communautaire dans une plateforme adapt\'ee aux groupes organis\'es.
